\chapter{\bilingual{GİRİŞ}{INTRODUCTION}}
\label{ch:introduction}

% Note: in the source, a single line break is just a space in the output, and a
% blank line starts a new paragraph. Writing one sentence per line (as below)
% keeps Git diffs small without changing the compiled result.
\bilingual{%
  Bu bölüm, tezin amacını, kapsamını ve katkılarını tanıtmaktadır.
  İstanbul Kültür Üniversitesi öğrencilerinin tezlerini resmi kılavuza uygun
  biçimde hazırlamaları için geliştirilen bu şablon, biçimsel ayrıntıları
  otomatikleştirir.}{%
  This chapter introduces the aim, scope, and contributions of the thesis.
  The template, developed so that Istanbul Kultur University students can
  prepare their theses in conformance with the official guide, automates the
  formatting details.}

\section{\bilingual{Motivasyon}{Motivation}}
\bilingual{%
  Biçimlendirme kurallarını elle uygulamak zaman alıcı ve hataya açıktır.
  Bu şablon, kenar boşluklarından kaynakça stiline kadar tüm kuralları tek bir
  yerde toplar.}{%
  Applying formatting rules by hand is time-consuming and error-prone.
  This template gathers every rule -- from margins to bibliography style -- in
  one place.}

\subsection{\bilingual{Amaç}{Objective}}
\bilingual{%
  Amaç, derlendiğinde kurallara tam uyumlu bir tez üreten, bakımı kolay bir
  LaTeX altyapısı sunmaktır.}{%
  The objective is to provide a maintainable LaTeX infrastructure that produces
  a fully compliant thesis when compiled.}

\section{\bilingual{Tezin Organizasyonu}{Organization of the Thesis}}
\bilingual{%
  Tezin geri kalanı şöyle düzenlenmiştir: ikinci bölüm ilgili çalışmaları,
  üçüncü bölüm yöntemi sunar, son bölüm ise sonuçları özetler.}{%
  The remainder of the thesis is organized as follows.
  Chapter~2 reviews related work, Chapter~3 presents the method, and the final
  chapter summarizes the conclusions.}
